\section{Extracting the modular aliases}
\label{sec:aliases}

The modular checksum includes subsets whose integer sum is not the target.
This section lists precisely those unwanted subsets, so that their phase
contributions can be subtracted.  We use multiple representations of a subset
by two masks on a small common block, together with modular filtering.  This
representation method and the distinction between many and few distinct
subset sums have precedents in
\cite{HGJ2010,AKKN2016,NederlofWegrzycki2021}.  The additional feature here is
an exclusion rule: pairs with the exact target weight are skipped without
being visited, even when their masks overlap.  That rule is what makes the
number of visited pairs small.

Throughout this section, fix the transformed weights $c_i$ and targets $t_s$
from \eqref{eq:transformed}.  Their bit lengths are $O(b+\log n)$.
The low-deficiency condition \eqref{eq:low-deficiency} will be used only
to prove coverage; the list and pair bounds hold without it.  The algorithm
does not test this condition.  It is the case in which the preliminary
exact procedure need not already decide the instance.  For the present
construction we impose the input-length guard
\begin{equation}
 b\le 2^{.00001n}.
 \label{eq:bit-guard}
\end{equation}
The final algorithm uses an exact fallback when this guard fails.

Choose a prime $P$ uniformly from
$[2^{\lfloor\eta n\rfloor},2^{1+\lfloor\eta n\rfloor})$, where
$\eta=.51772$, and put $b_i=c_i\bmod P\in\{0,\ldots,P-1\}$.
For a target $t_s$, define its set of \emph{aliases} by
\begin{equation}
 \mathcal A_s(P)=
 \left\{I\subseteq[n]:
     \sum_{i\in I}c_i\equiv t_s\pmod P,
     \quad \sum_{i\in I}c_i\ne t_s\right\}.
 \label{eq:alias-definition}
\end{equation}

\begin{theorem}[Alias extraction]
\label{thm:aliases}
There is one randomized procedure with the following properties for all
sufficiently large $n$.  Its input consists of the transformed integers
$c_i=Ca_i+u_i$ and targets $t_s=Ct+s$, $0\le s\le nN_0$, of
\eqref{eq:transformed}, with the input-length guard \eqref{eq:bit-guard}.
It chooses $P$ uniformly from
$[2^{\lfloor\eta n\rfloor},2^{1+\lfloor\eta n\rfloor})$
among the primes and returns either failure or $P$ together with a
deduplicated list of genuine aliases for each target.
If $\mathbb E_{|V|=\lfloor n/2\rfloor}d(V)\le .01004n$, then, with
probability $1-o(1)$ over $P$ and the remaining random choices, it does not
fail and its lists equal
$\mathcal A_s(P)$ simultaneously for every $s$.  The error bound is uniform
over the transformed inputs satisfying these assumptions.

For every outcome of its random choices, the procedure stores and visits at
most $2^{\lceil\tau n\rceil}\operatorname{poly}(n)$ records, where
$\tau=.48996$, and takes
$2^{\tau n}\operatorname{poly}(n+b+2)$ elementary bit operations, apart
from the implementation of exact uniform finite draws.  Those draws admit
the bounded implementations in Section~\ref{sec:implementation}.
\end{theorem}

We first construct the filtered lists and show why the exact-integer
exclusion keeps their enumeration small.  We then prove coverage: a
random prime preserves the distinctions needed for many representations
of each alias, and the smaller-modulus filter retains one of them with
sufficient probability.

\subsection{Representations and filtered lists}

Set
\begin{equation}
 m=2\left\lfloor\frac{.0567n}{2}\right\rfloor,
 \qquad \gamma=.48925.
 \label{eq:alias-parameters}
\end{equation}
Thus $m$ is even.  All arguments below concern sufficiently large $n$,
so $m>0$ and $17m/2\le\lfloor n/2\rfloor$.
Write $h(x)=-x\log x-(1-x)\log(1-x)$ for binary entropy, with
$h(0)=h(1)=0$.
The numerical bounds $h(.0567)+.0567<.40$ and $h(1/4)<.81128$ used
below follow, for example, from the logarithm series
\eqref{eq:certificate-log}, after scaling its argument into $[1,2]$.

Every alias $I$ of $t_s$ has a carry $q\in\{0,\ldots,n\}$ such that
\begin{equation}
 \sum_{i\in I}b_i=T_{s,q},
 \qquad T_{s,q}=(t_s\bmod P)+qP.
 \label{eq:alias-carry}
\end{equation}
For each $s$ and every $q=0,\ldots,n$, perform $n^{10}$ independent trials,
conditional on $P$ and the transformed input.  A trial chooses a uniform
$m$-element set $M\subseteq[n]$ and an independent uniform subset
$F\subseteq[n]$.  These two draws have the same law for every $P$ and
use randomness independent of it.
The block $M$ supplies multiple representations: selected positions in
it will be split between two masks.  For any fixed subset $I$, the flipped
set $J=I\mathbin\triangle F$ selects exactly half of this block with
inverse-polynomial probability, independently of the original cardinality
of $I$.
Introduce signed weights and translated targets
\begin{equation}
 \begin{aligned}
  b_i^F&=(-1)^{\mathbf1_{i\in F}}b_i,
  &T^F&=T_{s,q}-\sum_{i\in F}b_i,\\
  c_i^F&=(-1)^{\mathbf1_{i\in F}}c_i,
  &t_s^F&=t_s-\sum_{i\in F}c_i.
 \end{aligned}
 \label{eq:alias-flip}
\end{equation}
For every subset $I$, these definitions give
\[
 \sum_{i\in J}b_i^F=\sum_{i\in I}b_i-\sum_{i\in F}b_i,
 \qquad
 \sum_{i\in J}c_i^F=\sum_{i\in I}c_i-\sum_{i\in F}c_i.
\]
Thus flipping is a bijection, not a relaxation of the subset problem.

Within the trial, consider every $k=0,\ldots,m/2$.  Partition the positions
outside $M$ into ordinary left and right sets $L,R$, of sizes $l,r$,
where $l+r=n-m$.  Choose $l$ so that
\begin{equation}
 N_L=2^l\binom mk,
 \qquad N_R=2^r\binom m{m/2-k}
 \label{eq:alias-raw-sizes}
\end{equation}
are within a factor of two.  Such a choice follows by rounding
\[
 l=\frac{n-m+\log\binom m{m/2-k}-\log\binom mk}{2}
\]
to the nearest integer.  The unrounded value lies between
$(n-2m)/2$ and $n/2$, so the required partition exists.  The choice can
be made by exact comparisons of integers, without logarithmic arithmetic.
For definiteness, take the first $l$ positions of $[n]\setminus M$ as $L$.

A left record is a pair $(A,X)$ with $A\subseteq L$, $X\subseteq M$,
$|X|=k$; a right record is a pair $(B,Y)$ with $B\subseteq R$,
$Y\subseteq M$, $|Y|=m/2-k$.  Each record stores both its signed reduced
weight, obtained using $b_i^F$, and its signed unreduced weight, obtained
using $c_i^F$.  At this point $X$ and $Y$ need not be disjoint.
If they are disjoint, their union with $A,B$ represents a subset with
exactly $m/2$ selected positions in $M$.

For this $k$, independently draw a prime
\[
 P'\in[2^{\lfloor\gamma m\rfloor},2^{1+\lfloor\gamma m\rfloor})
\]
uniformly and then draw a uniform residue $j'\bmod P'$.  Keep left records
whose signed reduced weight is $j'\bmod P'$, and right records whose
signed reduced weight is $T^F-j'\bmod P'$.
Only these filtered records are formed explicitly.

To generate a filtered list, split its ordinary coordinate set into two
parts as evenly as possible.  Enumerate subsets of one part together with
all its permitted masks in $M$, and enumerate subsets of the other part.
Sorting these two lists by their weights modulo $P'$ permits their
matching pairs to be emitted in time polynomial in $n+b+2$ times the
sum of their lengths and the number of emitted records.  Indeed, for each
entry in one list, a binary search finds the interval of complementary
residues in the other, after which only output records are visited.
The setup work for both sides is at most
\begin{equation}
 \operatorname{poly}(n+b+2)
     2^{m+\lceil\max(l,r)/2\rceil}
 \le \operatorname{poly}(n+b+2)2^{.3067n+O(1)}.
 \label{eq:alias-generation}
\end{equation}
Here $\max(l,r)\le n/2+O(1)$ follows from the balancing rule.

\paragraph{The exact-integer exclusion.}
From the filtered lists, enumerate only pairs whose signed reduced
weights add to $T^F$ as integers and whose signed unreduced weights
do \emph{not} add to $t_s^F$ as integers.  Both sums count mask positions
with their multiplicities, including a position selected by both masks.
The excluded pairs need not be generated: group the right records by
signed reduced weight, and within each group sort by signed unreduced
weight.  For a left record, locate the group with the required reduced
weight.  The forbidden unreduced equality is a single contiguous interval
inside that group.  Two binary searches locate its endpoints, and the
algorithm visits only records outside that interval.

For each visited pair, test whether $X\cap Y=\varnothing$.  If so,
undo the flip on $A\cup B\cup X\cup Y$, check
\eqref{eq:alias-definition} directly, and save the resulting alias for
target $s$.  Deduplicate the saved subsets separately for each $s$.
In particular, every saved record is a genuine alias, whether or not
the procedure eventually finds all aliases.

\subsection{Why the visited lists and pairs are small}

The next estimates concern the uncapped construction.  They hold without
the low-deficiency assumption and average over the original random
choice of $P$.  In particular, we do not restrict $P$ to primes for
which the later coverage argument succeeds.

Conditional on all choices except $j'$, any fixed left or right record
survives with probability $1/P'$.  Balancing
\eqref{eq:alias-raw-sizes}, the entropy bound for binomial coefficients,
and concavity of $h$ give
\begin{align}
 \log\mathbb E[\text{one filtered list length}]
 &\le \frac{n-m}{2}
    +\frac m2\left(h(k/m)+h(1/2-k/m)\right)
    -\gamma m+O(1)\notag\\
 &\le \frac n2+\left(h(1/4)-\frac12-\gamma\right)m+O(1)
 <.48992n
 \label{eq:alias-list-bound}
\end{align}
for sufficiently large $n$.  The last inequality uses $h(1/4)<.81128$;
the coefficient obtained from this upper bound is $.489909101$.

To bound visited pairs, fix $s,q,k$, the block $M$, the flip $F$, and
two \emph{raw} records.  The indices $s,q,k$ range over predetermined
integers; in particular $q$ is fixed while $P$ is averaged.
These records and their ordinary-coordinate
partition can be fixed independently of $P$.  Let $\mu_i\in\{0,1,2\}$
be their combined multiplicity at position $i$, and define
\begin{equation}
 D=\sum_i(-1)^{\mathbf1_{i\in F}}\mu_i c_i
       -t_s+\sum_{i\in F}c_i.
 \label{eq:alias-fixed-difference}
\end{equation}
This is a fixed integer independent of $P$; in particular neither the
reduced weights nor the carry target occurs in its definition.
The exact-integer exclusion requires $D\ne0$.  The reduced-weight equality
implies $P\mid D$, because $T_{s,q}\equiv t_s\pmod P$ and
$b_i\equiv c_i\pmod P$.
The integer $D$ has $O(b+\log n)$ bits, even when masks overlap.
Lemma~\ref{lem:prime-collision} therefore gives
\begin{equation}
 \Pr_P[\text{both integer tests hold}]
 \le\operatorname{poly}(n)(b+\log n)2^{-\eta n}.
 \label{eq:alias-large-modulus-bound}
\end{equation}
For $D=0$ the pair is always excluded, so its contribution is zero.

Conditional on $P,P',M,F$ and the raw records, the smaller-modulus filters
can both hold for at most one value of $j'$.  Their probability is
therefore at most $1/P'\le 2^{1-\gamma m}$, regardless of whether the
reduced-weight equality holds.  Combining this bound with
\eqref{eq:alias-large-modulus-bound} and summing over the $N_LN_R$ raw
pairs yields
\begin{align}
 \mathbb E[\text{visited pair count}]
 &\le \operatorname{poly}(n)(b+\log n)
      2^{n+(2h(1/4)-1-\gamma)m-\eta n}\notag\\
 &<2^{.48988n}
 \label{eq:alias-pair-bound}
\end{align}
uniformly under \eqref{eq:bit-guard}, for sufficiently large $n$.
Indeed, using $h(1/4)<.81128$ and the guard, the coefficient of $n$
before polynomial factors is at most $.489848677$.

This is the reason to exclude exact equalities before testing mask
disjointness.  If exact equalities for overlapping pairs were visited,
they could have $D=0$ and would not satisfy the nonzero-divisor estimate.
The sorted-interval exclusion avoids their cost entirely.

Cap each filtered list and each visited-pair count at
\begin{equation}
 C_{\mathrm{alias}}=2^{\lceil\tau n\rceil}.
 \label{eq:alias-cap}
\end{equation}
The procedure returns failure before emitting a record that would exceed
a cap.  There are only polynomially many choices of $s,q$, trial, and
$k$.  Markov's inequality and
\eqref{eq:alias-list-bound}--\eqref{eq:alias-pair-bound} show that a cap
is exceeded with probability $o(1)$, uniformly over guarded inputs.
On every execution, including one that fails, list generation,
interval searches, mask checks, storage, and deduplication take at most
$2^{\tau n}\operatorname{poly}(n+b+2)$ bit operations, by
\eqref{eq:alias-generation} and the caps.

\subsection{Why every alias is recovered}

It remains to prove that the constructed representations cover every
alias.  The low-deficiency condition will provide many distinct mask
weights on a randomly selected block.  We first show that, with high
probability, reduction modulo $P$ preserves all distinctions of this kind.

\begin{lemma}[Sparse distinctions]
\label{lem:sparse-distinction}
With probability $1-o(1)$ over $P$, the following implication holds
simultaneously for every $v\in\{-1,0,1\}^n$ with at most $m$ nonzero
coordinates:
\begin{equation}
 \sum_i v_ic_i\ne0
 \quad\Longrightarrow\quad
 \sum_i v_ic_i\not\equiv0\pmod P.
 \label{eq:sparse-distinction}
\end{equation}
The error is uniform under \eqref{eq:bit-guard}.
\end{lemma}

\begin{proof}
There are at most
\[
 \sum_{j=0}^m\binom nj2^j
 \le 2^{(h(.0567)+.0567)n+O(\log n)}
 <2^{.40n}
\]
such vectors for sufficiently large $n$.  Every nonzero integer
$\sum_i v_ic_i$ has $O(b+\log n)$ bits.  Lemma~\ref{lem:prime-collision}
therefore bounds the probability that $P$ divides any one of them by
$\operatorname{poly}(n)(b+\log n)2^{-\eta n}$.
A union bound, followed by \eqref{eq:bit-guard}, gives a total error at most
$\operatorname{poly}(n)2^{-(\eta-.40-.00001)n}=o(1)$.
\end{proof}

In particular, fix any block of at most $m$ positions and any assignment
of signs on that block.  Distinct integer subset sums of the signed $c_i$
remain distinct modulo $P$ on the event of
Lemma~\ref{lem:sparse-distinction}.  Differences of two such sums have
coefficients in $\{-1,0,1\}$ on that block.  Consequently they also give
distinct \emph{integer} subset sums of the corresponding signed $b_i$.
The event covers all blocks and all sign assignments at once, including
those chosen by the random trials.

Assume now the low-deficiency condition \eqref{eq:low-deficiency}, and fix
a prime satisfying Lemma~\ref{lem:sparse-distinction}.  We continue in
the uncapped experiment.  Fix a target $s$ and an alias
$I\in\mathcal A_s(P)$, and use its carry $q$ in
\eqref{eq:alias-carry}.  The constants in the following probability bounds
do not depend on $I,s$, or the chosen prime.

\paragraph{A block with many distinct sums.}
In one trial, $J=I\mathbin\triangle F$ is a uniform subset of $[n]$
independent of $M$.  Hence
\begin{equation}
 \Pr[|J\cap M|=m/2]
   =2^{-m}\binom m{m/2}\ge\frac1{m+1}.
 \label{eq:alias-half-block}
\end{equation}
Conditional on this event, $S=J\cap M$ is uniform among the $m/2$-element
subsets of $[n]$.  To bound its expected deficiency, take a uniform
$\lfloor n/2\rfloor$-element set and partition it uniformly into
17 disjoint sets of size $m/2$ and a remainder.  Each of the 17 sets
has the same marginal law as $S$.  Superadditivity and nonnegativity of
deficiency, from Lemma~\ref{lem:deficiency}, therefore imply
\[
 17\mathbb E d(S)
 \le\mathbb E_{|V|=\lfloor n/2\rfloor}d(V)
 \le .01004n.
\]
Markov's inequality now gives, for all sufficiently large $n$,
\begin{equation}
 \Pr[d(S)\le .0106m\mid |J\cap M|=m/2]
 \ge 1-\frac{.01004n}{17(.0106)m}>.005.
 \label{eq:alias-useful-block}
\end{equation}
The limiting lower bound is greater than $.017$.

On this event the $c_i$ on $S$ have at least
$2^{(.5-.0106)m}=2^{.4894m}$ distinct subset sums.
For every realized $S,F$, changing signs translates their support:
\[
 \left\{\sum_{i\in X}c_i^F:X\subseteq S\right\}
 =\left\{\sum_{i\in X}c_i:X\subseteq S\right\}
       -\sum_{i\in F\cap S}c_i.
\]
The same number of distinctions survives for signed reduced integer
weights, by the fixed prime's sparse-distinction property.
Partitioning these subsets by cardinality, there is a
$k\in\{0,\ldots,m/2\}$ whose masks give at least
\begin{equation}
 B\ge\frac{2^{.4894m}}{m+1}
 \label{eq:alias-many-weights}
\end{equation}
distinct signed reduced weights.

For this $k$, fix the corresponding ordinary partition $L,R$ and combine
each such mask $X\subseteq S$ with $J\cap L$.  The resulting left records
still have $B$ distinct weights.  Every one has the complementary right
record $(J\cap R,S\setminus X)$; their masks have the prescribed sizes
and are disjoint.  Thus the representations all give the same flipped
alias $J$.  It remains to show that the independent smaller modulus
preserves enough of these choices to hit a random filter residue.

\paragraph{A random smaller modulus hits a representation.}
Let $w_1,\ldots,w_B$ be the distinct signed reduced weights just obtained.
They lie in $[-nP,nP]$, so every nonzero difference has $O(n)$ bits.
The choice of $k$ and these weights depends only on
$P,M,F,I$, not on the independent $P',j'$ used for that layer.
For a fixed $P'$, let
\[
 z_a=|\{j:w_j\equiv a\pmod {P'}\}|,
 \qquad a\in\mathbb Z/P'\mathbb Z.
\]
Applying Lemma~\ref{lem:prime-collision} to each nonzero difference gives
\begin{equation}
 \mathbb E_{P'}\sum_a z_a^2
 \le B+K_0 n^2B^2 2^{-\gamma m}
 \label{eq:alias-collision-moment}
\end{equation}
for an absolute constant $K_0$ and sufficiently large $n$.
Since $.4894>\gamma=.48925$, \eqref{eq:alias-many-weights} implies
$B\ge2^{\gamma m}$ for sufficiently large $n$.  The first term can
therefore be absorbed into the second by increasing $K_0$.
With probability at least $1/2$ over $P'$, Markov's inequality gives
\[
 \sum_a z_a^2\le 2K_0n^2B^2 2^{-\gamma m}.
\]
On this event, Cauchy--Schwarz shows that the number of occupied residues
is at least
\[
 \frac{(\sum_a z_a)^2}{\sum_a z_a^2}
 \ge \frac{2^{\gamma m}}{2K_0n^2}.
\]
As $P'<2^{1+\gamma m}$, a uniform $j'$ hits one of these residues
with probability at least $1/(4K_0n^2)$.

Whenever this happens, the corresponding left record survives, and its
complementary right record survives because their integer reduced weights
sum to $T^F$.  Their unreduced weights sum to
$\sum_{i\in I}c_i-\sum_{i\in F}c_i\ne t_s^F$, since $I$ is an alias.
Thus this pair is not excluded, is visited, passes the disjointness test,
and saves $I$.

Combining \eqref{eq:alias-half-block},
\eqref{eq:alias-useful-block}, and the smaller-modulus probability, a
single trial recovers $I$ with probability at least
\[
 \frac{.005}{m+1}\cdot\frac12\cdot\frac1{4K_0n^2}
 \ge n^{-4}
\]
for sufficiently large $n$.  The useful layer may depend on the alias
and on the trial: every layer is tried, with independent smaller-modulus
randomness, so no choice of that layer is required by the algorithm.

\paragraph{Simultaneous coverage and the caps.}
Conditional on $P$ and the transformed input, the $n^{10}$ trials for
the fixed $s,q$ are independent.  The probability of missing $I$ in all
of them is at most
$(1-n^{-4})^{n^{10}}\le e^{-n^6}$.
There are at most $(nN_0+1)2^n$ target--subset pairs.  A union bound
therefore shows that all aliases are recovered with conditional
probability $1-o(1)$.  This assertion holds for every prime satisfying
Lemma~\ref{lem:sparse-distinction}.

Finally couple the capped procedure to this uncapped experiment, using
the same random choices until a cap is exceeded.  The probability of
such an exceedance is $o(1)$ by the unconditional size estimates.  We do
not condition those estimates on a favorable prime or on successful
coverage.  Adding the sparse-distinction, coverage, and cap failure
probabilities proves the simultaneous conclusion of
Theorem~\ref{thm:aliases}.  Its every-execution operation bound was proved
above; bounded implementations of the finite draws add only the stated
implementation overhead and a further uniformly vanishing failure
probability.
